\renewcommand{\Anno}{1996}

\renewcommand{\NumeroProblema}{5}

\renewcommand{\Testo}{
La Scuola Superiore S. Anna custodisce con
estrema cura in una cassaforte gli elaborati dei concorsi di
ammissione.  Per motivi di sicurezza, la Direzione ha deciso di dotare
la cassaforte di un certo numero di serrature e $4$ impiegati
custodiscono ciascuno un certo numero di chiavi.  Per aprire la
cassaforte servono le chiavi di tutte le serrature, ma pi\`{u}
impiegati possono avere la chiave della stessa serratura.

Determinare il numero minimo di serrature di cui deve essere dotata la
cassaforte affinch\'{e}, distribuendo bene le chiavi, si abbia che 3
qualunque dipendenti siano in grado di aprirla, ma 2 qualunque
dipendenti non siano in grado di aprirla.  
}

\renewcommand{\Risposta}{6}

\renewcommand{\Aiutino}{
L'idea fondamentale \`{e} che ad ogni coppia di dipendenti deve 
mancare (almeno) una chiave. Pu\`{o} la stessa chiave mancare a 2 
coppie diverse di dipendenti? Quindi le chiavi devono essere almeno 
\ldots. A questo punto basta esibire un esempio in cui il numero 
di chiavi appena trovato \`{e} sufficiente.
}

\renewcommand{\SoluzioneUfficiale}{ Se 2 impiegati da soli non possono
aprire la cassaforte, ne segue che a loro manca almeno una chiave.  La
chiave di una stessa serratura non pu\`{o} mancare a 2 coppie diverse
di impiegati perch\'{e}, se cos\`{\i} fosse, unendo le 2 coppie
avremmo un gruppo di 3 o 4 impiegati (a seconda che le 2 coppie
abbiano o meno un elemento in comune) che non pu\`{o} aprire la
cassaforte, il che \`{e} contro l'ipotesi.

Ne segue che a coppie diverse mancano chiavi di serrature diverse, 
e quindi le serrature sono almeno tante quante le coppie di 
impiegati, e cio\`{e} 6.

Mostriamo ora che con 6 serrature (indicate con A, B, C, D, E, F) 
possiamo verificare le condizioni distribuendo le chiavi secondo lo 
schema indicato nella seguente tabella.
\begin{center}
	\renewcommand{\arraystretch}{1.3}
	\begin{tabular}{|c|c|}
		\hline
		Impiegato & Serrature che pu\`{o} aprire  \\
		\hline\hline
		1 & ABC  \\
		\hline
		2 & ADE  \\
		\hline
		3 & BDF  \\
		\hline
		4 & CEF  \\
		\hline
	\end{tabular}
\end{center}

\`{E} infatti immediato controllare che, comunque si scelgano 3
impiegati (si pu\`{o} fare in 4 modi), questi possiedono tutte le
chiavi, mentre comunque si scelgano 2 impiegati (si pu\`{o} fare in 6
modi), a questi manca sempre esattamente una chiave, come indicato
nella tabella seguente.
\begin{center}
	\renewcommand{\arraystretch}{1.3}
	\begin{tabular}{|c|c|}
		\hline
		Impiegati presenti & Chiave mancante  \\
		\hline\hline
		1--2 & F  \\
		\hline
		1--3 & E  \\
		\hline
		1--4 & D  \\
		\hline
		2--3 & C  \\
		\hline
		2--4 & B  \\
		\hline
		3--4 & A  \\
		\hline
	\end{tabular}
\end{center}
}

\renewcommand{\SoluzioniAlternative}{
}

\renewcommand{\Osservazioni}{\TitoloOsservazione Se le serrature sono 
6, la distribuzione delle chiavi tra gli impiegati \`{e} 
sostanzialmente forzata (a meno di permutazioni delle serrature stesse). 
Infatti, per la discussione fatta, sappiamo che in tal caso vi \`{e} 
una corrispondenza biunivoca tra coppie di impiegati e serrature. 
Scritta una qualunque di tali corrispondenze, come nella seconda 
tabella, risultano determinate le chiavi che ciascuno non possiede, 
dunque anche le chiavi che ciascuno possiede.}

\renewcommand{\Commenti}{
}